\chapter{Les fichiers : lequel contient quoi}

Vous allez accumuler des fichiers de tailles très différentes. Ce chapitre dit
lequel garder, lequel partager, et lequel effacer sans regret.

\section{Vue d'ensemble}

\begin{center}
\begin{longtable}{p{2.5cm}rp{9.0cm}}
\toprule
\textbf{Extension} & \textbf{Taille} & \textbf{Contenu} \\
\midrule
\endhead
\texttt{.gguf} & 4,4 Go & Le modèle de base, quantifié. Vient d'ailleurs, ne
   change jamais. \\
\texttt{.nkla} & 77 Mo & Les adaptateurs entraînés. \textbf{C'est ce qu'on
   partage.} \\
\texttt{\_ckpt0/1.nkla} & 231 Mo & Checkpoint de reprise. Personnel, temporaire. \\
\texttt{.nkgp} & variable & Checkpoint d'un GPT entraîné de zéro. \\
\bottomrule
\end{longtable}
\end{center}

\section{Le GGUF : le modèle de base}

C'est le gros fichier téléchargé au chapitre~2. Il contient les sept milliards de
paramètres, \textbf{quantifiés} — c'est-à-dire rangés sur moins de bits que leur
précision d'origine, pour tenir en mémoire.

\begin{center}
\begin{tabular}{lp{9.8cm}}
\toprule
\textbf{Q4\_K} & Environ 4 bits par nombre. La majorité des tenseurs. \\
\textbf{Q6\_K} & Environ 6 bits. Réservé aux tenseurs les plus sensibles — dont
   \texttt{output.weight}, celui qui choisit le mot final. \\
\textbf{F32} & 32 bits, pleine précision. Les normalisations, petites et
   critiques. \\
\bottomrule
\end{tabular}
\end{center}

\begin{nkretenir}
Sans quantification, ce modèle occuperait plus de 28~Go et ne tiendrait sur
aucune carte grand public. C'est elle qui rend tout ce cours possible.
\end{nkretenir}

\textbf{Ce fichier n'est jamais modifié par l'entraînement.} C'est le principe
même du LoRA : le socle est gelé. Vous pouvez le mettre en lecture seule.

\section{Le fichier .nkla : ce que vous avez appris}

\begin{center}
\begin{tabular}{lr}
\toprule
Taille & 77 Mo \\
Paramètres & 20\,185\,088 \\
Contenu & les pièces LoRA, et rien d'autre \\
\bottomrule
\end{tabular}
\end{center}

C'est le résultat de votre travail. Pour l'utiliser, il faut aussi le GGUF
d'origine : le \texttt{.nkla} seul ne sert à rien, comme une paire de lunettes
sans personne pour les porter.

\subsection{Ce que contient l'en-tête, et pourquoi}

Le fichier commence par sa carte d'identité : nombre de couches, de projections,
rang, alpha, et les dimensions du modèle auquel il est destiné.

Ces dimensions sont redondantes avec le GGUF — c'est voulu. Elles permettent de
\textbf{refuser} un adaptateur qu'on tenterait d'appliquer au mauvais modèle,
plutôt que de l'appliquer de travers et de produire un modèle silencieusement
faux.

\subsection{L'empreinte}

Le fichier se termine par une empreinte calculée sur tout son contenu. Au
chargement, elle est recalculée et comparée. Si elle diffère, le fichier est
\textbf{refusé}.

\begin{nkpiege}
Un adaptateur corrompu et chargé « au mieux » donnerait un modèle qui répond
encore — mais mal, de façon subtile. Vous chercheriez l'erreur dans votre corpus
pendant des jours. Un refus franc est toujours préférable à une dégradation
silencieuse.
\end{nkpiege}

\section{Les checkpoints : votre filet}

\begin{center}
\begin{tabular}{p{5.2cm}r}
\toprule
Les paramètres LoRA & 77 Mo \\
Le premier moment d'Adam & 77 Mo \\
Le second moment d'Adam & 77 Mo \\
L'état d'entraînement & quelques octets \\
\midrule
\textbf{Total} & \textbf{231 Mo} \\
\bottomrule
\end{tabular}
\end{center}

Les moments pèsent deux fois les paramètres : c'est toute la différence entre 77
et 231~Mo.

L'état d'entraînement, minuscule, est pourtant ce qui rend la reprise correcte :
numéro du pas, position dans le corpus, numéro d'époque, état du tirage
aléatoire, taux d'apprentissage et réglages d'Adam.

\begin{nkretenir}
Un checkpoint est un fichier \emph{personnel et temporaire}. On ne le partage
pas : votre interlocuteur n'a que faire de vos moments d'Adam. Il veut le
\texttt{.nkla} de 77~Mo.
\end{nkretenir}

\subsection{Pourquoi deux fichiers}

\texttt{\_ckpt0} et \texttt{\_ckpt1} sont écrits en alternance. À tout instant,
l'un des deux est une génération complète et saine. Si une coupure survient
pendant l'écriture, ou si un disque abîme un fichier, l'autre reste.

\begin{nkcode}[]{Le repli en action}
checkpoint mon-modele_ckpt0.nkla ignoré :
  empreinte FNV incorrecte — fichier corrompu
REPRISE : pas 10, exemple 10/100017, époque 0
\end{nkcode}

\section{Que garder, que jeter}

\begin{center}
\begin{tabular}{p{3.4cm}p{4.0cm}p{6.4cm}}
\toprule
\textbf{Fichier} & \textbf{Après l'entraînement} & \textbf{Pourquoi} \\
\midrule
\texttt{.gguf} & \textbf{Garder} & Nécessaire pour utiliser vos adaptateurs. \\
\texttt{.nkla} & \textbf{Garder, sauvegarder} & C'est votre travail. 77~Mo se
   copient partout. \\
\texttt{\_ckpt0/1} & Jeter si vous ne comptez pas continuer & 462~Mo à deux, et
   ils ne servent qu'à reprendre. \\
Le corpus & \textbf{Garder} & Pour réentraîner, corriger, augmenter. Souvent
   plus précieux que le modèle lui-même. \\
\bottomrule
\end{tabular}
\end{center}

\begin{nkpiege}
\textbf{Sauvegardez votre corpus avant votre modèle.} Un modèle perdu se
réentraîne en quelques jours de machine ; un corpus perdu, ce sont des semaines
de travail humain. C'est lui, le bien rare.
\end{nkpiege}

\section{Partager son travail}

Pour donner votre modèle spécialisé à quelqu'un, deux choses :

\begin{enumerate}
  \item Le \textbf{nom exact} du modèle de base (par exemple
        \texttt{qwen2.5:7b}) — il le téléchargera lui-même.
  \item Votre fichier \texttt{.nkla} de 77~Mo.
\end{enumerate}

Ajoutez-y, par honnêteté, le \textbf{rang} utilisé : sans lui, votre
interlocuteur devra deviner l'argument \texttt{-{}-rank}.

\begin{nkcode}[]{Ce qu'il tapera}
ollama pull qwen2.5:7b
NKQwen2Ask --gguf=<son_blob> --adapters=votre-modele.nkla --rank=8
\end{nkcode}

\section{Récapitulatif des commandes}

\begin{center}
\begin{longtable}{p{6.0cm}p{8.0cm}}
\toprule
\textbf{Ce que vous voulez} & \textbf{La commande} \\
\midrule
\endhead
Vérifier que tout marche &
  \texttt{NKQwen2Chat <modele>} \\
Mesurer avant d'entraîner &
  \texttt{NKQwen2Ask -{}-gguf=<m> -{}-question="…"} \\
Vérifier son corpus &
  \texttt{NKQwen2Train … -{}-steps=2 -{}-fresh} \\
Entraîner &
  \texttt{NKQwen2Train -{}-gguf=<m> -{}-corpus=<c> -{}-out=<o>} \\
Reprendre &
  la \emph{même} commande \\
Recommencer de zéro &
  la même commande $+$ \texttt{-{}-fresh} \\
Juger le résultat &
  \texttt{NKQwen2Ask -{}-gguf=<m> -{}-adapters=<o>.nkla} \\
Entraîner un GPT de zéro &
  \texttt{NK\_GPT\_CONFIG=<f>.cfg NKGptTrain} \\
\bottomrule
\end{longtable}
\end{center}

\section{Pour finir}

Vous savez maintenant préparer une machine, constituer un corpus, entraîner un
modèle pendant des jours sans craindre les coupures, spécialiser ce modèle sur
votre domaine, et juger honnêtement s'il s'est amélioré.

Le reste est une affaire de patience et de corpus. Ce sont les deux seules
choses qu'aucun outil ne fera à votre place.

\begin{nkretenir}
Trois idées à emporter :

\textbf{1.} Le corpus décide de tout. Les réglages ne rattrapent jamais de
mauvais textes.

\textbf{2.} La perte de contrôle est le seul juge fiable pendant
l'entraînement ; vos questions préparées d'avance sont le seul juge fiable
après.

\textbf{3.} Sur dix jours, une coupure est certaine. C'est pourquoi la reprise
n'est pas un confort mais la condition même du travail.
\end{nkretenir}

\begin{nkexo}
Le projet complet, à faire en une semaine.

\textbf{Jour 1.} Choisissez un domaine que vous connaissez et dont aucun modèle
ne peut rien savoir. Écrivez dix questions de référence, et mesurez l'avant.

\textbf{Jours 2 et 3.} Constituez un corpus de 300 à 500 paires depuis vos
sources réelles.

\textbf{Jour 4.} Vérifiez le corpus, lancez l'entraînement, laissez tourner.

\textbf{Jour 5.} Mesurez l'après sur les dix questions. Notez ce qui a changé et
ce qui n'a pas changé.

\textbf{Jours 6 et 7.} Corrigez ce que la comparaison a révélé — presque
toujours le corpus — et recommencez.

Ce cycle est le métier. Tout le reste n'est que de l'outillage.
\end{nkexo}
